\section{The Door and the Log}\label{sec:door}

The ledger is a door, some checks, a log, views computed from the log, and one
generic operation besides: at the read seam it applies a declared operator to
pre-boundary market data, carrying an old number into a new frame. This section
builds the door, the checks, the log, and the views from four primitives ---
wallet, unit, move, and transaction; the read-seam operation is the subject of
Section~\ref{sec:data}. Together these are the whole of what the ledger does.

\subsection{Wallets, units, and the atomic move}

A \emph{unit} is a thing that can be held --- a share, a currency, a futures
contract, a bespoke note. A \emph{wallet} is a named partition of position space,
and its entire state is one signed number per unit. A positive number is a long
position; a negative number is a short. There is no separate machinery for shorts
and no separate register of ownership standing beside the balance: economic
ownership \emph{is} the wallet balance. If a wallet holds $+2{,}000$ of a unit, it
owns 2{,}000 of that unit. For an outright holding, long or short, the scalar
balance is the whole story: there is nothing else to keep in step. The one
exception proves the rule. A securities loan carries coordinates the balance alone
does not hold --- the borrow's recall risk, its fee accrual, the manufactured
dividend owed on the lent line --- and those are not smuggled into the balance;
they are added as their own units and moves at the securities-lending stop
(Section~\ref{sec:tour-sbl}). An outright long or short has no such coordinates,
and for it the balance says everything.

The only way a balance changes is a \emph{move}.

\begin{definition}[Atomic move]\label{def:move}
A move transfers a positive quantity $q$ of one unit from a source wallet to a
destination wallet in one indivisible step: the source decreases by $q$, the
destination increases by $q$. The magnitude is positive; the direction is carried
by which wallet is the source and which the destination.
\end{definition}

Two design decisions are folded into this definition, and both pay off later.
First, the quantity is an exact integer in minor units --- never a floating-point
number. A share count is an integer; a cash amount is an integer number of cents.
Division does occur --- a withholding on an odd amount, cash in lieu of a
fractional entitlement, an odd per-share dividend, a three-for-two on an odd lot
--- but it never leaves an unbooked remainder. The rounding convention is itself
declared data, and whatever a division cannot distribute evenly is booked
explicitly, as a cash-in-lieu or rounding leg that is a move like any other. The
conserved quantities stay exact not because nothing divides but because the
remainder is never dropped: it is recorded. Conservation and replay are therefore
matters of integer arithmetic, decided exactly at any scale. Second, a move carries a single positive magnitude and two wallets. It
has no field in which a quantity could be created or destroyed. Just as a move
has no endpoint beyond the two wallets it connects, the door admits no outcome
beyond success and refusal. Its Haskell signature says exactly that:

\begin{lstlisting}
applyTx :: Transaction -> Ledger -> Either LedgerError Ledger
\end{lstlisting}

A transaction either produces a whole new ledger or a typed \texttt{LedgerError}.
There is no third outcome and, in particular, no partially applied state.

A \emph{transaction} is the atomic event recorded in the log: an ordered list of
moves together with the non-balance changes that ride with them --- bookkeeping
rows, status writes, an introduction or appending of terms. It applies in full or
not at all. Atomicity is the property that pays off here: a corporate action that
must double a share count, halve a strike, and set a basis marker does all three
in one transaction or none of them, so no reader can ever observe a
half-applied event.

\subsection{Conservation, by construction}

Here is the law the whole design is organised around: for every transaction and
every unit, the total quantity across all wallets is unchanged. Nothing is created
and nothing is destroyed; value only moves.

The important thing is \emph{how} this law holds. It is not checked at the door.
Look again at what a move does: it subtracts $q$ from one wallet and adds $q$ to
another, so its net contribution to the total quantity of its unit is $q - q = 0$.
A transaction is a list of moves, and the total change it makes is the sum of the
per-move contributions, each of which is zero. The sum is zero because every term
is zero, not because anything counted it and found it balanced. There is no
transaction that fails to conserve waiting to be rejected --- such a transaction
cannot be built, because the only tool for changing a balance is a move, and a
move is a transfer.

This is worth dwelling on because it is the pattern the design repeats
everywhere. One could imagine a \texttt{set\_balance} operation that writes a
wallet's balance directly, guarded by a runtime check that total value still sums
to zero. That design has a hole: the check can be bypassed, disabled, or wrong,
and the day it is, the books are unconserved and nobody knows. The design here has
no \texttt{set\_balance}. Opening a position is not a special primitive; it is an
ordinary move from a \emph{virtual wallet} that stands for the outside
counterparty. Every counterparty outside the system is represented by a virtual
wallet inside it, so every move has two endpoints inside the ledger and the system
is closed. A conservation failure is therefore not a rejected financial event; it
is an impossible one --- and if it ever appeared it would be an implementation
bug, not a transaction the ledger admitted.

\subsection{The single door and what it checks}

Every change to the ledger goes through one function. Registration goes through
it; a trade goes through it; a corporate action goes through it. There is no
second way in, no privileged channel, and no writer of a balance other than a move
applied by this door. This matters precisely because contracts are untrusted: a
contract cannot hold write access, because write access does not exist as a thing
to hold. A contract can only submit a proposal, and a proposal becomes fact only
by passing through the same door as everything else.

Since conservation cannot fail, the door does not check it. The door checks the
things that \emph{can} fail:

\begin{itemize}[leftmargin=1.5em]
\item \textbf{Basis agreement.} A quantity and a price may be combined only when
they are expressed in the same frame. When a transaction records or consumes an
observation, the door checks that its basis marker agrees with the state it is
being joined to. This is the check that a stale, pre-split price is not silently
multiplied against a post-split share count; Section~\ref{sec:data} explains the
frame it enforces.
\item \textbf{Idempotence.} Each transaction carries an identifier derived from
its cause. If the same transaction is submitted twice --- a processor retried
after a crash, a timer that fired late and then fired again --- the door commits
it once. A coupon due on one date is paid once even if its contract runs three
times.
\item \textbf{Legal status writes.} Facts other than balances --- a unit's
lifecycle state, its basis marker, its last settlement --- are written only
through a small closed set of status writes, each applied by a single writer as an
absolute replacement. A contract cannot invent a new kind of status write or
write a field that is not its to write.
\item \textbf{The invariance weld.} A transaction that carries a unit across a
basis boundary is admitted only if its booked entitlement moves realise the
boundary's declared operator: for the boundary's $(f, c)$, each holder's quantity
leg realises $q \mapsto q \cdot f$ and each cash leg equals $q \cdot c$. This welds
an adjustment's moves to the operator the boundary declares, so a transaction
cannot claim a boundary and then book moves that do not match it. Later sections
apply the same weld in two forms derived from this one --- across two boundaries
sharing an effective time, and per holder or in aggregate when holders elect ---
but the check is this one, stated here.
\end{itemize}

Alongside these, the door refuses a transaction that references a unit never
registered, that points its cause at an event not yet committed, or that would
leave the adjustment schedule incomplete --- each a well-formedness condition,
each returning a typed error and no state. The common thread is that the door
never guesses and never repairs. A rejected transaction leaves no trace in the
ledger's state and no entry in the log: \texttt{applyTx} returns either the whole
new ledger or a typed error, and a typed error is a returned value, not a commit.
The refusal carries the rejected payload back to whoever submitted it; at the datum
ingest door this returned refusal \emph{is} the quarantine, holding the payload as
a value without ever storing it inside the system of record. If a deployment keeps
a durable journal of what it refused, that journal is an operational artifact on
the submitting side, never a second log within the ledger. Defaults fail closed:
an absent declaration is a refusal, not an assumed identity; an unregistered unit
is refused with its payload returned, not admitted on optimistic terms. The door would rather stop than improvise, because improvising
is exactly how a system of record acquires a fact nobody can account for.

\subsection{The log, and why untrusted proposals are safe}

What the door admits, it appends to a log that is never rewritten. The log is
ordered and hash-chained: no earlier entry is edited in place, and any tampering
with history is evident to any party holding a prior chain head --- anchoring that
head outside the system is a deployment matter. Corrections are not edits; a mistake is
repaired by a later, compensating transaction that is itself recorded, so the
record of the mistake and the record of its repair both survive.

We can now answer the question the thesis raises: if contracts are untrusted, why
is it safe to record what they propose? Because admission and recomputation
together bound what a bad proposal can do. A defective contract cannot break an
invariant --- there is no unconserved proposal for the door to admit, no partial
state it can leave, no row it can delete. The worst it can do is book wrong
economics inside a well-formed transaction, and because everything it read is in
the log, that error is not hidden: anyone can recompute the transaction from the
log and find it. This detection is after the fact. The recomputation runs over the
committed log, and nothing gates a view or the settlement projection on its having
run first; the guarantee is that a wrong economics is always findable from the log
and detected when the witness runs --- which the deployment must schedule --- not
verified before it is consumed. What that buys is precise. The structural invariants are never
breached --- not for an instant, because they hold by construction, not by a check
a wrong contract might slip past. The error is confined to the \emph{economics} of
the one well-formed transaction that carried it, and its repair is deterministic,
with full lineage, through the compensating path. So the blast radius of a wrong
contract is bounded above by ``a detectable, well-formed, wrong transaction,
repaired by a recorded correction'', which is a bounded and repairable thing. It is
never ``a corrupted ledger'', which is not. The gates that do act before
consumption --- the confirmation gate that will not let a datum cross an unarmed
boundary (a boundary is armed only once its operator is confirmed), the pending
quarantine on a datum with no operator to carry it across ---
guard boundary-\emph{crossing} consumption, not a producer's economics; they are
not recomputation and must not be mistaken for it.

Consider the first transaction any book records: a purchase of 1{,}000 shares of a
company at \EUR{100}. It enters as one transaction with one move --- 1{,}000 of
the share unit from the counterparty's virtual wallet into the book's wallet ---
and a second move of \EUR{100{,}000} of cash the other way. The share count in the
book's wallet is now $1{,}000$; the counterparty's virtual wallet is $-1{,}000$,
which is the reconciliation point against the counterparty's own record of the
same trade. Nothing here is a stored ``position'' that some later process must
keep current. The position is $1{,}000$ because the log says one move of
$1{,}000$ happened, and it will say so forever.

\subsection{Every view is a fold of the log}

A balance is the sum of the moves that touched a wallet. A valuation over a chosen
set of wallets is the sum of their balances times an externally supplied price
vector; profit and loss is the difference of two valuations. That difference
depends only on the endpoints, not on the route between them --- but it measures
performance only when no external contribution crossed into the wallet set between
the two dates. A deposit or a withdrawal moves the valuation without being profit;
performance is the endpoint difference \emph{less} the net external flow over the
interval. Inside the closed system every such flow is itself a recorded move across
the wallet set's boundary, so the adjustment is exact and computed from the log,
never estimated. With flows accounted for, the endpoints alone fix the number.
Value, unlike quantity, obeys no conservation law: a book can be worth more today
than yesterday, and nothing is violated. What matters is that every one of these
numbers is \emph{computed} from the log when asked, not stored beside it. The
unit-status cache the ledger keeps for speed is exactly that --- a cache, derivable
by replaying the events and safe to discard, never an authority. Replaying a
unit's events to any past point rebuilds the exact value that held then.

This is what makes the reconciliation failures of
Section~\ref{sec:thesis} unrepresentable rather than rare. The balance sheet and
the profit-and-loss statement are two folds of one log. For them to disagree,
one fold would have to see an event the other does not, and there is only one log.

\subsection{Time travel re-reads; it never re-executes}

Because the log is complete and the views are folds of it, the ledger can
reconstruct its exact state as of any past moment: replay the stream up to that
moment and stop. Balances, unit state, and the basis coordinate all come back
exactly as they stood. This reconstruction re-reads recorded transactions; it
never re-executes a contract, and it fetches nothing from outside the log. The
consequence is worth stating precisely.

This is the property that makes an upgradeable, untrusted contract layer safe to
live with: \emph{a contract can be upgraded, replaced, or discarded, and every
past state of the ledger still reconstructs exactly}, because reconstruction reads
the transactions a contract once produced and never re-runs the contract that
produced them. A contract's output, once admitted, is a fact in the log like any
other. Whether the contract still exists, still behaves the same way, or was
thrown away years ago has no bearing on what the log says happened. Replay is
deterministic and independent of where it is cut, so the same history
reconstructs to the same state every time. The record does not depend on the
machinery that fed it --- which is the property a system of record must have, and
the reason the ledger can hold untrusted output without being at the mercy of it.
